\documentclass[Total.tex]{subfiles}

\begin{document}
	\chapter{Preparations}
		\section{Class, and its Operations}
		\section{Order Relations}
			\begin{Def}
				Consider the following first-order language
				\begin{itemize}
					\item 
					\begin{gather*}
						\mathcal{L}_{ord,strict}=(M,<)
					\end{gather*}
					
					with: $M$ is a set, $<\subseteq M\times M$, and 
					\begin{itemize}
						\item Irreflexivity
						\begin{gather*}
							\forall x \neg(x<x)
						\end{gather*}
						\item Transitivity
						\begin{gather*}
							\forall x\forall y\forall z((x<y\wedge y<z)\rightarrow x<z)
						\end{gather*}
					\end{itemize}
					
					Thus, 
					\begin{gather*}
						T_{so}=\begin{Bmatrix}
							\forall x \neg(x<x)\\
							\forall x\forall y\forall z((x<y\wedge y<z)\rightarrow x<z)
						\end{Bmatrix}\qquad \mathcal{M}\models T_{so}\Leftrightarrow \mathcal{M}\mbox{ is strictly partial order set}
					\end{gather*}
					
					\item 
					\begin{gather*}
						\mathcal{L}_{ord}=(M,\{\leq,=\})
					\end{gather*}
					
					with $M$ is a set, $\leq \quad\subseteq M\times M$, 
					\begin{itemize}
						
						\item Reflecxitive
						\begin{gather*}
							\forall x(x\leq x)
						\end{gather*}
						\item Anti-Symmetric
						\begin{gather*}
							\forall x\forall y(x\leq y\wedge y\leq x\rightarrow x=y)
						\end{gather*}
						\item Transitivity
						\begin{gather*}
							\forall x\forall y\forall z((x\leq y\wedge y\leq z)\rightarrow x\leq z)
						\end{gather*}
					\end{itemize} 
					
					Thus,
					\begin{gather*}
						T_{o}=\begin{Bmatrix}
							\forall x \neg(x<x)\\
							\forall x\forall y(x\leq y\wedge y\leq x\rightarrow x=y)\\
							\forall x\forall y\forall z((x\leq y\wedge y\leq z)\rightarrow x\leq z)
						\end{Bmatrix}\qquad \mathcal{M}\models  T_\mathcal{o}\Leftrightarrow \mathcal{M}\mbox{ is a partial order}
					\end{gather*}
					
				\end{itemize}
			\end{Def}
			\subsection{Linear-Ordering}
			
			\begin{Def}
				Especially, a partial order $(M,<)$ is called \textbf{linear ordered}, if
				\begin{gather*}
					\forall x\forall y(x<y\vee y<x\vee x=y)
				\end{gather*}
				
				Thus,
				\begin{gather*}
					\mathcal{T}_{lo}=\begin{Bmatrix}
						\forall x \neg(x<x)\\
						\forall x\forall y\forall z((x<y\wedge y<z)\rightarrow x<z)\\
						\forall x\forall y(x<y\vee y<x\vee x=y)
					\end{Bmatrix}\qquad \mathcal{M}\models T_{lo}\Leftrightarrow \mathcal{M}\mbox{ is linear ordered}
				\end{gather*}


			\end{Def}
			\subsubsection{Well-Ordering}
			
			\begin{Def}
				Let $(X,\leq)$ be linearly ordered. The relation $\leq$ is called \textbf{well-ordering} if every nonempty subset of $X$ has a least element. 
				\begin{gather*}
					\forall A\subseteq X,A\not=\emptyset\Rightarrow \exists a\in A\forall b\in A,a\leq b
				\end{gather*}
				
				Thus,
				\begin{gather*}
					\mathcal{T}_{wo}=\begin{pmatrix}
						\forall x \neg(x<x)\\
						\forall x\forall y(x\leq y\wedge y\leq x\rightarrow x=y)\\
						\forall x\forall y\forall z((x\leq y\wedge y\leq z)\rightarrow x\leq z)\\
						\forall A\subseteq X,A\not=\emptyset\Rightarrow \exists a\in A\forall b\in A,a\leq b
					\end{pmatrix}\qquad \mathcal{M}\models  \mathcal{T}\Leftrightarrow \mathcal{M}\mbox{ is well-ordered}
				\end{gather*}
			\end{Def}
		\subsection{Well-Founded Order}
			\begin{Def}
				The well-founded order is the generation of well-ordering
				\begin{gather*}
					\forall A\subseteq X,A\not=\emptyset\Rightarrow \exists a\in A\forall b\in A,a\leq b
				\end{gather*}
			\end{Def}
			
		\subsection{Density}
			\begin{Def}
				内容...
			\end{Def}
	\chapter{ZFC Axioms}
		\section{ZFC Axioms}
			
			\subsection{ZFC Axioms}
			
			
			\subsection{Class, Properclass}
		\section{Axiom of Regularity}
			\begin{proposition}
				There exists no set $x\in x$. 
			\end{proposition}
			\begin{proof}
				内容...
			\end{proof}
			
			\begin{corollary}
				There is no pair of sets $x,y$ such that
				\begin{gather*}
					x\in y\qquad y\in x
				\end{gather*}
			\end{corollary}
			\begin{proof}
				内容...
			\end{proof}
			
			\subsection{Von Neumann Hierarchy, Rank}
				\begin{Def}
					内容...
				\end{Def}
		\section{Axiom of Choice}
			
			\begin{proposition}
				Assume: ZF axiom, the following statements are equivalent:
				\begin{itemize}
					\item The axiom of choice;
					\item (Choice Function)If $x$ is such that $y\not=\emptyset,\forall y\in x$, then there exists a function
					\begin{gather*}
						f:x\rightarrow \cup x\qquad f(y)\in y
					\end{gather*}
					\item (Zorn's Lemma) Let $(x,\leq)$ be a partial order set. Then, suppose
					\begin{itemize}
						\item $x$ is nonempty;
						\item Every chain has a upper bound. i.e. 
						\begin{gather*}
							\forall y(y\subseteq x\wedge (a\in y\wedge b\in y\rightarrow a\subseteq b\vee b\subseteq a))
						\end{gather*}
					\end{itemize}
					
					Then, $P$ has at least one maximal element $x$ such that: $\forall a(a\in x\rightarrow x\not\subset a)$. 
					\item (Hausdorff's Maximal Principle) Let: $X\subseteq P(X)$ be a set, for $Y\subseteq X$, with 
					\begin{gather*}
						\forall a\forall b(a\in Y\wedge b\in Y\rightarrow a\subseteq b\vee b\subseteq a)
					\end{gather*}
					
					Then, there must some $a\in X$ such that
					\begin{gather*}
						\forall b(b\in Y\rightarrow b\subseteq a)
					\end{gather*}
					
					That is, some $a\in X$ such that there is no $b\in X$, with $a\subset b$. 
					\item (well-ordered theorem)Every set can be well-ordered
					
					\item (Cartesian Product) The Cartesian product of any family of nonempty sets is nonempty. 
				\end{itemize}
			\end{proposition}
			\begin{proof}
				\begin{itemize}
					\item (2)$\Leftrightarrow$ (3)
					\item (2)$\Leftrightarrow$ (4)
					\item (2)$\Leftrightarrow$ (5)
					\item (2)$\Leftrightarrow$ (6)
				\end{itemize}
			\end{proof}
	\chapter{Construction of Real Numbers}
		\section{Natural Numbers}
		\section{Integer Numbers}
		\section{Rational Numbers}
		\section{Real Numbers}
			\begin{axiom}
				We could construct real number $(\RR,+,\times,<) as following$:
				
				\begin{itemize}
					\item Field Axioms:
					\item Order Field Axioms:
					\item Archimede Field: 
				\end{itemize}
			\end{axiom}
			\subsection{Compeleteness Condition}
				
				\begin{proposition}
					The followings are equivalent:
					\begin{itemize}
						\item Metric space $(\RR,|-|$) is compelete;
						\item For any upper subset $M\subseteq \RR$, $sup(M)$ exists;
					\end{itemize}
				\end{proposition}
				
				\subsubsection{Cauchy Sequence}
				
				\begin{lemma}
					Suppose $M\subseteq \RR$ be upper bounded, then $sup(M)$ exists. 
				\end{lemma}
				\begin{proof}
					Select upper bound $N$ and $x\in M$,
					
					\begin{gather*}
						a_0:=x,b_0=N,c_n:=\frac{a_n+b_n}{2}\quad
						a_{n+1}:=\begin{cases}
							a_n&\forall_{x\in M} c_{n}\geq x\\
							c_{n}&Elsewise
						\end{cases}\quad b_{n+1}=\begin{cases}
							b_n&\exists_{x\in M} c_n<x\\
							c_n&Elsewise 
						\end{cases}
					\end{gather*}
					
					$(a_n),(b_n)$ are monotone.
					
					Then, $|a_n-b_n|\leq 2^{-(n-1)}|a_1-b_1|\Rightarrow|b_{n+m}-b_n|\leq |b_n-a_n|\leq 2^{-(n-1)} |b_1-a_1|$. It is Cauchy sequence. Thus $b_n\rightarrow z\in \RR$. And $a_n\rightarrow z$: 
					
					$z$ is $sup(M)$. Suppose $y\in M$, we have
					\begin{gather*}
						y>z\Leftrightarrow \exists N_0(n>N_0\rightarrow y>a_n)
					\end{gather*}
					
					Contradict to the definition.
				\end{proof}
				
				\begin{lemma}
					If we define $\RR$ as
					\begin{gather*}
						\RR:=2^\NN
					\end{gather*}
					
					Then, 
				\end{lemma}
					
		\chapter{Forcing}
		
			Our preparations:
			\begin{itemize}
				\item Models, Languages;
				\item Filters;
			\end{itemize}
			
			\subsection{Transitive Model}
				\begin{Def}
					Recall: Traisitive set;
					\begin{itemize}
						\item 
						A set $M$ is called \textbf{transitive} if
						\begin{gather*}
							\forall x\in M,x\subseteq M
						\end{gather*}
						
						equivalently, $\cup M\subseteq M$;
						
						\item Consider $\mathcal{L}_\in=\{\in\}$, with $\mathcal{M}=(M,E_\in)$, then, $M$ is called a \textbf{traisitive model} of theory $T$ if
						\begin{itemize}
							\item $M$ is a transitive set;
							\item $E=\in |_{M\times M}$;
							\item $(M,\in)\models T$, where $T$ is a set-theoretic theory;
						\end{itemize}
					\end{itemize}
				\end{Def}
			
			\section{Forcing Notion}
				\begin{Def}
					Take the following notion $\mathbb{P}$ for the \textbf{forcing notion}, with
					\begin{gather*}
						\mathbb{P}=(P,\leq,1_\mathbb{P})
					\end{gather*}
					\begin{itemize}
						\item $M^\mathbb{P}$ by recusion:
						\begin{gather*}
							M_0^\mathbb{P}:=\emptyset\\
							M_{\alpha+1}^\mathbb{P}=P^M(\mathbb{P}\times M_\alpha^\mathbb{P})\\
							M_\alpha^\mathbb{P}=\cup_{\beta<\alpha} M_\beta^\mathbb{P}\\
							M^\mathbb{P}=\cup_{\alpha\in ord} M_\alpha^{\mathbb{P}}
						\end{gather*}
						
						which is called \textbf{$\mathbb{P}$-Names}
						\item $(P,\leq)$ is a preorder;
					\end{itemize}
					
					
					Here, $M$ is the universe of set theory. 
				\end{Def}
			
\end{document}